% ============================================================
%  chapters/ch14-accessibility-landscapes.tex
% ============================================================

\chapter{Accessibility Landscapes}

\chapterepigraph{%
  A landscape has peaks, valleys, and paths between them.\\
  Intelligence is knowing which paths lead where.
}{Flyxion, \textit{Semantic Infrastructure}}

\sidenote{App.\,J gives the\\full formal\\treatment of\\accessibility\\geometry}

Part IV begins the formal development of constraint theory and admissibility.
This chapter introduces accessibility landscapes — the geometric representation
of future possibility — as the central object of the framework.

\section{Admissible Futures}

From a given state $x$ at time $t$, some future states are reachable
and some are not. Among those reachable, some are admissible (they satisfy
all constraints) and some are not. The accessibility landscape is the
function that maps each present state to the volume of its admissible futures.

\begin{definition}{Admissible Future Set}{adm-future-set}
  The \emph{admissible future set} from state $x$ over horizon $H$ is:
  \[
    \mathcal{A}_H(x) = \{y \in X : \exists\,\text{admissible path }
    \gamma : x \rightsquigarrow y \text{ with } |\gamma| \leq H\}
  \]
  where admissibility of the path requires $\kappa(\gamma(t)) = 0$ for
  all $t \in [0, |\gamma|]$.
\end{definition}

\begin{definition}{Accessibility Landscape}{access-landscape}
  The \emph{accessibility landscape} is the function
  $S_H : X \to \mathbb{R}_{\geq 0}$ defined by:
  \[
    S_H(x) = \log \mathrm{Vol}\bigl(\mathcal{A}_H(x)\bigr).
  \]
  Points with high $S_H$ have many admissible futures; points with low $S_H$
  are near dead ends. The landscape $S_H$ is a scalar field over the
  state space.
\end{definition}

\section{Constraint Geometry}

The shape of the accessibility landscape is determined entirely by the
constraint structure. Tight constraints create narrow valleys; loose
constraints create broad plateaus.

\begin{definition}{Constraint Tightness}{constraint-tightness}
  The \emph{tightness} of a constraint $c$ at state $x$ is:
  \[
    \tau_c(x) = \frac{\|\nabla_x \kappa_c(x)\|}{\kappa_c(x) + \epsilon}
  \]
  for small $\epsilon > 0$. Large tightness indicates the constraint is
  actively binding (small $\kappa_c$, large gradient).
\end{definition}

\begin{proposition}{Tightness Reduces Accessibility}{tightness-access}
  For any state $x$ and active constraint $c$, increasing the tightness
  $\tau_c(x)$ (holding all else fixed) decreases $S_H(x)$.
\end{proposition}

\begin{proof}[Sketch]
  An active constraint eliminates the directions in its gradient direction
  from the admissible future set. Higher tightness means the constraint
  is more binding — it eliminates more of the local tangent space — reducing
  the volume of admissible futures.
\end{proof}

\section{Navigation Over Possibility Spaces}

The accessibility landscape provides a natural framework for \emph{navigation}:
moving through state space in ways that preserve or increase future options.

\begin{definition}{Accessibility-Preserving Path}{acc-preserving-path}
  A path $\gamma : [0,T] \to X$ is \emph{accessibility-preserving} if
  $S_H(\gamma(t))$ is non-decreasing in $t$.
\end{definition}

Accessibility-preserving paths move through state space in ways that
keep future options open. They avoid dead ends, narrow canyons, and
irreversible commitments. They represent \emph{flexible} strategies —
strategies that maintain the ability to respond to new information.

\begin{definition}{Accessibility Gradient Flow}{access-gradient-flow}
  The \emph{accessibility gradient flow} is the dynamical system:
  \[
    \dot{x}(t) = \nabla_x S_H(x(t)).
  \]
  Trajectories of this flow always move in the direction of increasing
  accessibility — toward states with more admissible futures.
\end{definition}

\begin{remark}
  Accessibility gradient flow is not always the best strategy. Sometimes
  the optimal path requires temporarily reducing accessibility (entering
  a narrow canyon) to reach a high-accessibility plateau on the other side.
  This is the mathematical version of the common observation that the
  most flexible long-term strategies sometimes require short-term commitments.
\end{remark}

\section{Accessibility Wells, Peaks, and Ridges}

The landscape $S_H$ has topological features that correspond to important
structural properties of the underlying system.

\begin{definition}{Landscape Features}{landscape-features}
  \begin{itemize}
    \item An \emph{accessibility well} is a local minimum of $S_H$:
          a state with fewer admissible futures than all nearby states.
          Wells are traps: systems that enter them tend to stay.
    \item An \emph{accessibility peak} is a local maximum of $S_H$:
          a state with more admissible futures than all nearby states.
          Peaks are hubs: systems near them have the most flexibility.
    \item An \emph{accessibility ridge} is a connected set of high-$S_H$
          states: a corridor of flexibility connecting accessible regions.
          Ridges are the highways of the possibility space.
  \end{itemize}
\end{definition}

\begin{example}
  In a career trajectory:
  \begin{itemize}
    \item A highly specialized technical skill is an accessibility well:
          expertise is deep but narrow; few adjacent career paths exist.
    \item A broad education in mathematics, writing, and programming is
          an accessibility peak: many careers are open; the system
          has maximum flexibility.
    \item A sequence of jobs in related but broadening fields is a ridge:
          each step maintains accessibility while moving toward
          a destination.
  \end{itemize}
\end{example}

\begin{exercise}
  Define the \emph{escape time} from an accessibility well as the
  expected number of steps needed to reach a state with $S_H$ above
  a threshold $\tau$, starting from the well minimum. How does escape
  time depend on the depth of the well (how much lower $S_H$ is at the
  minimum than at the surrounding threshold)?
\end{exercise}

\begin{exercise}
  Show that the gradient flow $\dot{x} = \nabla S_H(x)$ cannot escape
  from an accessibility well (cannot reach a state with higher $S_H$
  than the well entrance). What does this say about the limitations of
  pure accessibility-gradient strategies?
\end{exercise}

\begin{exercise}[title={Cosmological}]
  In the RSVP framework, structures (galaxies, stars, organisms) are
  interpreted as accessibility wells in the cosmological accessibility
  landscape. Explain why such structures are stable: once a galaxy forms,
  why does it persist? Use the accessibility well interpretation rather
  than conventional gravitational arguments.
\end{exercise}
